\chapter{Calcul littéral et identités}
\label{corriges:05}
\section*{Exercice 5.1 -- Développer et réduire}
$A=6x^2+15x-4x-10-4x=6x^2+7x-10$.
Pour $x=-1$, $A=6-7-10=-11$. Vérification dans la forme initiale : $(-5)\times3+4=-11$.

\section*{Exercice 5.2 -- Factoriser}
\[
12x^2-8x=4x(3x-2),\quad25x^2-9=(5x-3)(5x+3),
\]
\[
(x+1)(2x-3)+4(x+1)=(x+1)(2x+1).
\]
Par exemple $(5x-3)(5x+3)=25x^2+15x-15x-9=25x^2-9$.

\section*{Exercice 5.3 -- Deux programmes}
Le programme A donne $(x+4)^2=x^2+8x+16$, exactement l'expression du programme B. Cette identité établit l'égalité pour tout réel $x$. Pour $x=-3$, le résultat vaut $(-3+4)^2=1$ (ou $9-24+16=1$).

\section*{Exercice 5.4 -- Une propriété d'entiers}
$(n+2)^2-n^2=n^2+4n+4-n^2=4(n+1)$, donc c'est un multiple de $4$. Ensuite $101^2=(100+1)^2=10000+200+1=10201$. Les deux questions utilisent une identité valable pour tous les nombres, pas seulement une vérification numérique.
